% Part: model-theory
% Chapter: interpolation
% Section: introduction

\documentclass[../../../include/open-logic-section]{subfiles}

\begin{document}

\olfileid{mod}{int}{int}

\olsection{परिचय}

अंतर्वेशन प्रमेय म्हणजे पुढील निष्कर्ष: समजा
$\Entails !A \lif !B$. मग असे !!a{sentence}~$!C$ असते की
$\Entails !A \lif !C$ आणि $\Entails !C \lif !B$. शिवाय,
$!C$ मधील प्रत्येक !!{constant}, !!{function} आणि !!{predicate}
($\eq$ खेरीज) हे $!A$ आणि~$!B$ या दोन्हींमध्ये आढळते. !!{sentence}~$!C$
ला $!A$ आणि~$!B$ यांचे \emph{अंतर्वेशक} म्हणतात.

अंतर्वेशन प्रमेय स्वतःही रंजक आहे; पण त्याचे मुख्य महत्त्व असे की,
उपपत्तीतील परिभाष्यतेविषयीचे निष्कर्ष आणि दोन सुसंगत उपपत्ती एकत्र केल्यावर
मिळणारी उपपत्ती कोणत्या अटींवर सुसंगत असते याविषयीचे निष्कर्ष सिद्ध करण्यासाठी
ते वापरता येते. पहिला निष्कर्ष बेथचे परिभाष्यता प्रमेय म्हणून, तर
दुसरा रॉबिन्सनचे संयुक्त सुसंगतता प्रमेय म्हणून ओळखला जातो.

\end{document}
